\documentclass[a4paper,12pt]{ctexart}
\usepackage{xcolor}
\usepackage{graphicx}
\usepackage[top=1.8cm, bottom=1.8cm, left=1.8cm, right=1.8cm]{geometry}
\usepackage{array}
\usepackage{hyperref}
\usepackage{booktabs}
\hypersetup{colorlinks=true,linkcolor=blue!70!black}
\linespread{1.35}

\title{\textcolor{blue!70!black}{\Huge\textbf{金融学}}}
\subtitle{\textcolor{gray}{经济学门类 · 学生版}}
\author{}
\date{}

\begin{document}
\maketitle
\thispagestyle{empty}

\section{一句话概括}

金融学研究钱怎么流动、怎么管理、怎么增值。从个人理财到企业融资，从股票投资到风险控制，都在它的范围内。

\section{自我评估——你适合吗？}

\begin{itemize}
  \item[\textbf{☐}] 数学还不错，尤其概率统计不觉得烦
  \item[\textbf{☐}] 平时会留意"这个东西多少钱""为什么涨价了"
  \item[\textbf{☐}] 抗压能力还行，考试前不会崩
  \item[\textbf{☐}] 逻辑清楚，喜欢分析问题
  \item[\textbf{☐}] 不排斥和人打交道
  \item[\textbf{☐}] 细心，不容易算错数
\end{itemize}

\textbf{勾了4个以上}：这个专业值得认真考虑。
\textbf{勾了2-3个}：可以试试，但要想清楚自己哪里需要补。
\textbf{勾了1个以下}：建议看看别的门类。

\section{大学四年学什么}

\begin{tabular}{ll}
\toprule
大一 & 微观经济学、宏观经济学、高等数学、会计学 \\
大二 & 金融学原理、公司金融、统计学、计量经济学 \\
大三 & 投资学、金融工程、国际金融、财务报表分析 \\
大四 & 金融风险管理、实习、毕业论文 \\
\bottomrule
\end{tabular}

\section{北京高考选科要求}

\begin{tabular}{ll}
\toprule
院校 & 选科要求 \\
\midrule
北大光华/清华经管 & 必须选物理 \\
人大金融/央财 & 不限（建议选物理） \\
对外经贸/首经贸 & 不限 \\
\bottomrule
\end{tabular}

\section{这个专业最大的几个坑}

\begin{enumerate}
  \item 竞争极其激烈——金融分数线常年最高，同学都很强
  \item 头部吃肉尾部喝汤——最好的岗位一年只招几十人
  \item 学历很重要——本科进好机构越来越难，基本都要读研
\end{enumerate}

\end{document}
